\chapter{Affordance Fields and Cue Coherence}

\section{The Ladder and the Goldfish, Revisited}

Chapter 7 settled a question of principle with a deliberately simple example. A ladder affords climbing to a human and does not afford climbing to a goldfish, and since the ladder itself is unchanged between these two cases, whatever affordance is, it cannot be a property stored at the ladder's location in the field. It has to be a relation between the field and something the field does not itself contain: a capability, brought by whoever or whatever encounters it.

That example was chosen for its clarity, and clarity came at a cost. A ladder either affords climbing or it does not, for a given kind of visitor, and the binary nature of the example made the relational argument easy to see but easy to underestimate. Most affordances are not binary, and most encounters are not between a fixed object and a fixed species. Consider a toy car left on a table. To a small child, it affords pushing, racing, crashing into things, narrating a story about where it is going. To an adult tidying the room, it affords picking up and putting away. To a diagnostic engineer, if the car happens to be a scale model with functioning parts, it might afford inspection of a suspension mechanism. None of these are wrong. None of them exhaust what the car is. What the car affords is not one fact waiting to be discovered but a structured space of possibilities, different slices of which become available depending on who or what is doing the encountering, and depending on what that encounter is for.

This chapter develops affordance as exactly this: a structured, capability-relative space of possible actions, arising at the point where the persistent field meets an agent's capacity to act on it. The relation stated informally in Chapter 7,

\[ A : M \times C \to \mathcal{P}(\Delta\psi), \]

is this chapter's central object. M is the field context surrounding the point of encounter; C is the capability or cognitive state of the agent doing the encountering; the output is a set of viable proposals, candidates for the kind of change, \(\Delta\psi\), that agent might now put forward. The rest of this chapter is concerned with what actually populates that output set, and with how it comes to be populated in the first place.

\section{Cues: What Triggers an Affordance}

Affordances do not present themselves all at once. A courtyard, once observed and settled into appearance as Chapter 9 described, does not hand its visitor an exhaustive list of everything that could ever be done there. Something has to draw attention to a particular possibility before it becomes live for that agent at that moment, and this book will call that something a cue.

A cue is a localized signal, arising from the field, that activates relevance for a particular capability. To avoid a collision with this book's own notation, since \(R_t\) already names the field's causal residue, cue-triggered relevance will be written \(\rho_c(x)\) rather than reusing R: a scalar field over the relevant space, peaking near whatever has drawn the cue's attention and falling away with distance from it. A bench's silhouette is a cue that activates sitting-relevance for a capability profile that includes the ability to sit. A crack running along the base of the wall is a cue that activates inspection-relevance for a capability profile that includes structural expertise, and activates nothing in particular for a capability profile that does not. The cue does not manufacture the affordance out of nothing; it draws a specific region of the field's already-existing structure into the foreground of a specific agent's attention. Behavior, in this framework, follows the gradient of activated relevance: an agent moves toward and engages with whatever has been drawn into the foreground, in proportion to how strongly it has been drawn there.

\section{How Affordances Are Learned}

A single encounter establishes that a given cue was, on that occasion, associated with a given affordance. It does not yet establish a stable disposition. That stability accumulates the way most stable dispositions in this book's framework accumulate: through repetition, weighted and reinforced over time rather than fixed after one instance.

Write \(w_{ij}\) for the strength of the association between cue i and affordance j, for a given capability profile. Each encounter updates this weight by a simple rule,

\[ w^{t+1}_{ij} = (1-\alpha)\, w^t_{ij} + \alpha \cdot \mathrm{affordance}_{ij}, \]

where \(\mathrm{affordance}_{ij}\) is the strength with which that particular encounter realized affordance j in response to cue i, and α governs how quickly new encounters overwrite the accumulated history of old ones. This is nothing more exotic than a running average, but its consequences are worth spelling out. A cue that reliably predicts the same affordance across many encounters accumulates a strong, stable weight, and comes to trigger that affordance quickly and confidently on future encounters. A cue whose associated affordance varies from encounter to encounter never accumulates a strong weight toward any one of them, and continues to require deliberation each time it is met. This is how the same silhouette that once required careful inspection to identify as a bench, affording sitting only after some thought, comes eventually to trigger sitting-relevance immediately and without deliberation: not because the object has changed, but because the association between that cue and that affordance has been reinforced past the point of needing to be re-derived.

\section{The Vocabulary Affordances Draw From}

The affordance relation's output, the set P(\(\Delta\psi\)) that A(M, C) returns, is not an unstructured bag of arbitrary possibilities. It draws from a small, recurring vocabulary of operation types, and naming that vocabulary now will save considerable trouble later, when Part III develops what happens to a \(\Delta\psi\) once it has been proposed.

Five operations recur across nearly every affordance this book will need to discuss: store, retrieve, modify, compose, and rewrite. To store is to commit some content to a location for later retrieval, the operation underlying anything like note-taking, saving, or the invention-to-memory principle from Chapter 9 itself. To retrieve is to recover previously stored content, the operation underlying recall, lookup, and reference. To modify is to alter existing content in place, the operation underlying editing, repair, and revision. To compose is to combine multiple pieces of content into a new structure, the operation underlying assembly, construction, and synthesis. To rewrite is to replace content according to some rule rather than an arbitrary edit, the operation underlying transformation, translation, and structured revision. Climbing a ladder is a physical instance of composition, joining an agent's trajectory to the ladder's structure to reach a location otherwise unreachable. Sitting on a bench is closer to a temporary modification of the agent's own state. Inspecting a crack is retrieval, pulling latent structural information into the open. The vocabulary is small enough to be worth memorizing and general enough to cover far more than its origin, in a formalism developed originally for externalized memory systems rather than physical affordances, might suggest. This book draws on that vocabulary only for its operational alphabet; the fuller apparatus surrounding it, including results establishing that systems built from these five operations are computationally universal, belongs to that formalism's own treatment and will not be redeveloped here.

\section{Affordance Fields and Cue Coherence}

An agent rarely encounters a single cue in isolation. A courtyard offers a bench, a crack, a doorway, and a patch of shade all at once, each activating its own relevance, each potentially inconsistent with the others in what it recommends. Sitting and inspecting the crack may compete for the same attention; moving toward the doorway and lingering in the shade may pull in opposite directions. Something has to determine whether the affordances activated across these different cues cohere into a single usable policy, or fragment into contradictory pulls that leave the agent's next action undetermined.

This is a local-to-global consistency question, of exactly the shape a sheaf is built to answer, and it deserves to be stated as one now rather than left as a loose metaphor. Let CueCat be the category whose objects are cues or contexts and whose morphisms record how one cue's context relates to or refines another's. The affordance structure activated at each cue defines a presheaf,

\[ F_{\mathrm{cue}} : \mathrm{CueCat}^{\mathrm{op}} \to \mathrm{Set}, \]

assigning to each cue the set of affordances it locally activates, and to each relation between cues a way of restricting one context's affordances to a more specific one. The coherence question this chapter has been building toward is exactly the sheaf condition applied to \(F_{\mathrm{cue}}\): do the locally activated affordances, across overlapping cues and contexts, glue into a single section, a coherent policy the agent can actually act on, or do they fail to agree where their contexts overlap, leaving no consistent way to act on all of them at once?

It is worth being precise about what this sheaf is and is not a claim about, because a very different sheaf-theoretic construction is waiting in Chapter 21, and confusing the two would blur a distinction this book has gone to some trouble to keep sharp. \(F_{\mathrm{cue}}\)'s gluing question is epistemic and agent-relative: it asks whether one agent's locally triggered possibilities for action can be assembled into one coherent policy for that agent, given that agent's capability profile. It says nothing, on its own, about whether the world itself is coherent, whether the field's locally valid states assemble into one globally admissible reality. That is a separate question, concerning a separate sheaf, defined over a separate base category, and Chapter 21 will develop it as its own construction rather than as a restatement of this one. The two constructions are connected, not identical: a cue, after all, arises from an observation of the field, and Chapter 21's bridge section will make precise how observation maps cue-level coherence onto field-level admissibility. For now, this chapter's task is only to establish the first of the two sheaves cleanly, on its own terms, without borrowing the second one's vocabulary before it has been earned.

\section{Looking Ahead}

This chapter has given affordance a home appropriate to what it actually is: not a property stored in the field alongside coherence, flow, entropy, and appearance, but a relation, triggered by cues, learned through repetition, drawn from a small operational vocabulary, and coherent or incoherent across an agent's encounters in a sense now made precise. Something remains conspicuously absent from the account so far. Every affordance discussed in this chapter has been treated as though it arose fresh at the moment of encounter, with no reference to what has happened here before. But the courtyard has a history: the bench was placed at some point, the crack appeared at some point, and a full account of what a region affords ought to be able to draw on that history rather than reconstructing everything from the field's current appearance alone. That history is what \(R_t\), the causal residue introduced and then set aside in Chapter 7, was always meant to hold, and Chapter 11 is where this book finally gives it the treatment its role in the five-component decomposition has been waiting for.
